\documentclass[conference]{IEEEtran}
\usepackage[utf8]{inputenc} 
\usepackage[pdftex]{graphicx}
\usepackage[cmex10]{amsmath}
\usepackage{array}
\usepackage[hidelinks]{hyperref}
\usepackage[table,xcdraw]{xcolor}
\usepackage{booktabs}
\usepackage{multirow}
\usepackage{amsmath}
\usepackage{algorithm}
\usepackage[noend]{algpseudocode}
%\usepackage{subfig} 
\usepackage{amsfonts}
\usepackage[numbers]{natbib}
\usepackage{txfonts}
\usepackage{xcolor}

\DeclareMathOperator*{\argmax}{arg\,max}
\DeclareMathOperator*{\argmin}{arg\,min}

\newcolumntype{L}[1]{>{\raggedright\let\newline\\\arraybackslash\hspace{0pt}}m{#1}}
\newcolumntype{C}[1]{>{\centering\let\newline\\\arraybackslash\hspace{0pt}}m{#1}}
\newcolumntype{R}[1]{>{\raggedleft\let\newline\\\arraybackslash\hspace{0pt}}m{#1}}
\newcommand{\specialcell}[2][c]{\begin{tabular}[#1]{@{}c@{}}#2\end{tabular}}

\definecolor{azul}{HTML}{08519C}
\definecolor{vermelho}{HTML}{B22222}
\definecolor{verde}{HTML}{008000}

\newcommand{\AZ}{\textcolor{azul}{\blacksquare}}
\newcommand{\VER}{\textcolor{vermelho}{\blacksquare}}
\newcommand{\VERDE}{\textcolor{verde}{\blacksquare}}

\newcommand{\TG}{\textcolor{verde}{\blacktriangle}}
\newcommand{\TR}{\textcolor{vermelho}{\blacktriangledown}}


\makeatletter
\def\BState{\State\hskip-\ALG@thistlm}
\makeatother

\algnewcommand\algorithmicinput{\textbf{Input:}}
\algnewcommand\INPUT{\item[\algorithmicinput]}
\algnewcommand\algorithmicoutput{\textbf{Output:}}
\algnewcommand\OUTPUT{\item[\algorithmicoutput]}

\makeatletter
\newcommand{\linebreakand}{%
\end{@IEEEauthorhalign}
\hfill\mbox{}\par
\mbox{}\hfill\begin{@IEEEauthorhalign}
}
\makeatother

 
\begin{document}
 
\title{Mapeamento da Propagação de Eventos e Ajuste Residual em Modelos de Séries Temporais}
% \IEEEoverridecommandlockouts


\author{
 \IEEEauthorblockN{Jefferson dos Anjos}
 \IEEEauthorblockA{CEFET/RJ \\ jefferson.anjos@cefet-rj.br}
 \and
  \IEEEauthorblockN{Luiz Nogaroli}
 \IEEEauthorblockA{}
 \and
 \IEEEauthorblockN{Fernando Queiroz }
 \IEEEauthorblockA{}
 \and
 \IEEEauthorblockN{Felipe da Rocha Henriques}
 \IEEEauthorblockA{CEFET/RJ \\ felipe.henriques@cefet-rj.br}
 \and
 \IEEEauthorblockN{Diego Carvalho}
 \IEEEauthorblockA{CEFET/RJ \\ dcarvalho@ieee.org}
 \and
 \IEEEauthorblockN{Eduardo Ogasawara}
 \IEEEauthorblockA{CEFET/RJ \\ eogasawara@ieee.org}%
 % \thanks{Author’s Accepted Manuscript. Released under the Creative Commons license: Attribution 4.0 International (CC BY 4.0) \newline \protect\url{https://creativecommons.org/licenses/by/4.0/}}%
 }
	
%\IEEEpeerreviewmaketitle
\maketitle

\begin{abstract}
This paper proposes a hybrid method for short-horizon financial forecasting that combines quantitative models with semantic retrieval of comparable historical events. The pipeline detects events through return thresholds, structures news into JSON records, associates each event with its realized post-event sequence, and reuses this event library under walk-forward evaluation. The method retrieves semantically similar past events, applies a similarity-weighted impact adjustment, and learns an event-aware residual correction. Experiments on PETR4, PRIO3 and EXXO34, synchronized with Brent, generate 12,159 out-of-sample forecasts across LSTM, recurrent autoencoder and Transformer baselines. The hybrid method reduces average RMSE from 1.6248 to 0.8470 (47.9\%), reduces average MAE from 1.2445 to 0.6187 (50.3\%), and outperforms the base models on 76.8\% of the evaluated days. These results show that retrieving comparable historical events is a useful and auditable complement to purely quantitative forecasting.
\end{abstract}

\begin{keywords}
Keywords: Financial time series forecasting; Event retrieval; News-aware forecasting; Hybrid models; Semantic similarity; Residual correction.
\end{keywords}

\section{Introdução}

Mercados financeiros reagem continuamente a informações corporativas, macroeconômicas e geopolíticas. Embora essa premissa seja bem aceita, a incorporação operacional de noticias em modelos preditivos ainda permanece incompleta. Em boa parte da literatura, os modelos são treinados apenas com histórico de preços. Quando texto e incorporado, seu uso costuma ocorrer como uma variável contemporânea de sentimento ou como um canal adicional de classificação de movimento \cite{Tetlock:2007, Tetlock:2008, Xu:2018, Sawhney:2020}. Esse desenho deixa em aberto uma questão metodológica importante: como reutilizar eventos históricos semanticamente semelhantes e suas trajetórias observadas de impacto para melhorar previsões numéricas de curto prazo?

O problema e particularmente relevante em ativos expostos a choques recorrentes de commodities, regulação e geopolítica, como os do setor de óleo e gás. Nesses casos, o mesmo tipo de notícia tende a reaparecer em contextos diferentes, com efeitos que nem sempre se esgotam no dia do evento (que se propagam por dias). A abordagem proposta investiga essa lacuna ao converter notícias em eventos estruturados (investiga a propagação dos efeitos dessas notícias no preço do ativo e estima os efeitos quando um evento semanticamente equivalente volta a acontecer), estimar suas sequências realizadas de impacto e reaproveitar essas sequencias quando um evento semanticamente comparável volta a ocorrer.

A abordagem proposta organiza a camada informacional de forma modular. A avaliação offline principal recupera diretamente pares \textit{motivo-sequência} da biblioteca histórica de eventos, aplica um ajuste por impacto semântico com corte temporal estrito e depois aprende (ou apreende) uma correção residual em regime \textit{walk-forward}. Os clusters semânticos continuam existindo, mas cumprem papel principalmente exploratório e de auditabilidade.

Nossa hipótese de pesquisa (investiga se a recuperação de eventos semanticamente equivalentes e sua propagação no tempo produzem efeitos que, quando considerados, sejam capazes de reduzir o erro dos modelos base que costumam desprezar os efeitos das notícias em suas predições) é que recuperar eventos semanticamente comparáveis e suas sequencias observadas de impacto reduz o erro dos modelos base e melhora a auditabilidade da previsão em relação a abordagens puramente quantitativas.

Com base nisso, as contribuições centrais do artigo são:

\begin{itemize}
    \item uma biblioteca de eventos financeiros estruturados a partir de noticias, com propagação observada de impacto apos o evento (um repositório estruturado com notícias que impactam o ativo financeiro de interesse e uma estrutura de índices que representa o impacto da notícia nos dias subsequentes à data do evento);
    \item um mecanismo de recuperação semântica de eventos históricos comparáveis com restrição temporal explícita (?);
    \item um modelo híbrido em duas etapas, composto por ajuste sequencial por similaridade e por correção residual \textit{walk-forward};
    \item uma avaliação comparativa sobre três arquiteturas de base e três ativos do setor de óleo e gás, com evidencias quantitativas, ablações e logs de decisão.
\end{itemize}

\section{Referencial Teórico}

\subsection{Previsão financeira com modelos quantitativos}

Modelos baseados em aprendizado profundo têm sido usados com frequência na previsão financeira por sua capacidade de capturar relações não lineares e dependências temporais complexas. Redes LSTM mostraram desempenho competitivo em relação a abordagens tradicionais \cite{Fischer:2018}, enquanto Autoencoders recorrentes ajudam a construir representações latentes mais robustas a ruido \cite{Bao:2017}. Mais recentemente, arquiteturas baseadas em atenção, como o Transformer e suas adaptações para séries temporais, passaram a disputar esse mesmo espaço por sua eficiência em sequências longas e em cenários multi-horizonte \cite{Vaswani:2017, Lim:2021, Wu:2021}.

\subsection{Notícias, texto e eventos no mercado financeiro}

Há ampla evidência de que texto financeiro e fluxo de notícias afetam retorno, volume e percepção de risco \cite{Tetlock:2007, Tetlock:2008, Loughran:2011, Mitra:2011}. O desenvolvimento de embeddings distribuídos e de modelos contextualizados ampliou a capacidade de representar semelhança semântica entre frases e eventos \cite{Mikolov:2013, Devlin:2019, Araci:2019, Reimers:2019}. Ainda assim, grande parte dos trabalhos usa o texto como um sinal agregado no instante corrente, e não como um mecanismo explicito de recuperação de padrões históricos de impacto.

\subsection{Reação gradual e memória de eventos}

Embora a Hipótese dos Mercados Eficientes (HME) proponha a incorporação rápida da informação aos preços \cite{Fama:1970}, a literatura também documenta reações graduais, exageros e reversões \cite{Cutler:1989, Barberis:1998, Hong:1999}. Essa linha teórica abre espaço para metodologias que representem o (a propagação do) impacto de um evento como uma sequência temporal observável, e não apenas como um choque instantâneo. E precisamente nesse ponto que a recuperação semântica de eventos históricos se torna relevante: ela (, pois) oferece uma forma concreta de consultar experiências anteriores comparáveis e reutilizar suas trajetórias observadas (e estimar a propagação do seu impacto no tempo).

\section{Trabalhos Relacionados}

Os trabalhos mais próximos desta proposta (desta investigação) podem ser agrupados em quatro famílias. A primeira reúne modelos puramente quantitativos baseados em séries de mercado, como LSTM, Autoencoders e Transformers adaptados para previsão temporal \cite{Fischer:2018, Bao:2017, Lim:2021}. Essas abordagens servem como base de comparação natural, mas tratam choques informacionais apenas como parte da dinâmica endógena da série (mas ignoram o impacto das notícias em seu modus oeprandi).

A segunda família usa noticias, eventos ou textos de redes sociais para prever a direção do mercado \cite{Ding:2014, Ding:2015, Xu:2018, Sawhney:2020}. O ganho principal desses trabalhos está em incorporar (o sentimento apurado no texto) texto ao processo preditivo, mas, em geral, o problema é formulado como classificação de movimento, e o texto é fundido ao modelo ponta a ponta, sem gerar uma biblioteca auditável de sequências históricas reutilizáveis.

Além da literatura acadêmica, há implementações públicas amplamente divulgadas que combinam aprendizado profundo, múltiplas fontes de atributos e sinais textuais para previsão de mercado. Um exemplo é o repositório \textit{stockpredictionai}, de \cite{Banushev:2019}, que organiza uma proposta aplicada com GANs, LSTMs, CNNs e diferentes grupos de \textit{features}. Essa linha é relevante como referência prática, mas difere da proposta aqui apresentada (desse trabalho) por tratar o texto predominantemente como entrada do modelo preditivo, e não como uma biblioteca explícita (um repositório) de eventos históricos com sequências observadas e recuperação auditável.

A terceira família avança para arquiteturas com memória e interações mais estruturadas entre texto, tempo e múltiplos ativos \cite{Luo:2023}. Já a quarta família aproxima-se de esquemas de recuperação e memória baseados em LLMs \cite{Xiao:2025FinSrag, Wang:2025StockMem}. Em ambos os casos, a literatura recente oferece mecanismos sofisticados de memória, mas tende a depender de arquiteturas mais pesadas, foco em classificação ou uso dominante do modelo de linguagem como \textit{backbone}.

Em contraste, o presente artigo propõe uma alternativa mais leve e orientada à previsão numérica: a unidade principal de memória não é um estado latente do modelo, mas uma biblioteca explicita (um repositório) de eventos históricos com sequências observadas de impacto (e atribuição de fatores multiplicativos para seus impactos nos dias subsequentes) seus impactos; a recuperação é feita com corte temporal estrito; e a integração com o preditor ocorre por ajuste sequencial e correção residual. A lacuna preenchida, portanto, esta na combinação entre eventos estruturados, recuperação semântica direta, reutilização de sequências observadas e avaliação auditável sobre baselines quantitativos padrão.

\section{Metodologia}

\subsection{Escopo experimental e dados}

O escopo empírico do estudo abrange os ativos PETR4, PRIO3 e EXXO34, sincronizados com o Brent (BZ=F) como proxy de commodity. O pipeline trabalha com aproximadamente (recupera cerca de) cinco anos de dados diários provenientes do Yahoo Finance, cobrindo o período que, na base atual de eventos, vai de 10 de novembro de 2020 a 25 de março de 2026. O foco setorial foi deliberado: ativos de óleo e gás são fortemente sensíveis (impactados) a choques de notícias, (vasriações nos preços das commodities) commodity, (eventos geopolíticos) geopolítica e regulação, o que favorece a observação de efeitos recorrentes.

\begin{algorithm}[H]
\caption{Resumo do pipeline hibrido proposto}
\label{alg:pipeline}
\small
\begin{algorithmic}[1]
\STATE Detectar dias de evento por limiar de retorno absoluto
\STATE Estruturar cada evento em JSON com contexto textual e motivos
\STATE Construir a biblioteca histórica com pares \textit{motivo-sequencia} e datas
\STATE Gerar previsões base com LSTM, Autoencoder recorrente e Transformer
\FOR{cada data $t$ na avaliação \textit{walk-forward}}
    \STATE Recuperar eventos anteriores a $t$ com ativo compatível e similaridade semântica acima do limiar
    \STATE Agregar as sequencias recuperadas por media ponderada por similaridade
    \STATE Ajustar a previsão base no horizonte D0--D4
    \STATE Estimar a correção residual usando apenas a janela passada observada
\ENDFOR
\STATE Registrar previsões, eventos recuperados e metadados de auditoria
\end{algorithmic}
\end{algorithm}

\subsection{Detecção e estruturação de eventos}

Eventos são disparados quando a variação percentual diária absoluta do ativo ultrapassa um limiar predefinido (A busca e captura da notícia é iniciada quando a variação percentual diária ultrapassa o limiar estabelecido). (Nesse estudo, ) A versão principal do experimento usa um limiar de 2,0\%, mas a abordagem também gera analise de sensibilidade para 1,5\% e 2,5\%. Na configuração de 2,0\%, a base atual contem 320 eventos para PETR4, 469 para PRIO3 e 301 para EXXO34.

Quando um evento e detectado, o pipeline consulta um modelo da OpenAI (um modelo de LLM) para estruturar um registro JSON contendo data, ativo, retorno do dia, descrição do contexto, motivos identificados, sentimento e fontes. Quando o LLM não encontra contexto suficiente, o pipeline utiliza Tavily como \textit{fallback} de busca e reexecuta a estruturação do evento sobre o material recuperado. O conjunto experimental armazena atualmente 1.619 arquivos JSON de evento, com 3.253 motivos textuais identificados.

\subsection{Sequências realizadas e biblioteca de eventos (repositório de notícias)}

Cada evento recebe uma sequência realizada de impacto após sua ocorrência. O processo calcula retornos do dia do evento e dos pregões seguintes, armazenando sequências D0--D5 quando houver dados suficientes. Se um novo evento ocorre antes do fim da janela, a propagação e interrompida, evitando sobreposição ingênua de impactos. Ao considerar apenas registros com sequência não vazia, a biblioteca corrente de recuperação passa a conter 3.139 pares \textit{motivo-sequência}.

Formalmente, definimos a biblioteca histórica como
\[
\mathcal{L} = \{(m_i, q_i, d_i, a_i)\}_{i=1}^{N},
\]
em que $m_i$ e o motivo textual do evento, $q_i$ e a sequência observada de impacto, $d_i$ e a data do evento e $a_i$ e o ativo associado. Na inferência, apenas itens com $d_i < t$ e ativo compatível com a previsão corrente podem ser usados, o que reduz vazamento de informação futura.

\subsection{Clusters semânticos como camada exploratória}
Paralelamente a biblioteca direta de eventos (à construção do repositório de notícias), a abordagem (o método) constrói uma camada de clusterização semântica. Os motivos são embedados, podem ser (podendo ser) ampliados por variações textuais quando há poucos exemplos e são agrupados por \textit{Agglomerative Clustering}. Na base atual (Atualmente), essa etapa gera 200 clusters exploratórios, 50 para cada um dos quatro grupos (PETR4, PRIO3, EXXO34 e Brent). Esses clusters são uteis para auditoria qualitativa, para a pergunta de pesquisa sobre recorrência de padrões e para o MVP online. Na avaliação offline principal, entretanto, a unidade de recuperação e o evento histórico individual, e não o cluster agregado (achei esse trecho confuso e não consegui melhorar).

\subsection{Modelos base e ajuste hibrido}

Três modelos quantitativos foram treinados como baselines: LSTM, Autoencoder recorrente e Transformer. Todos recebem séries sincronizadas entre ativo (o ativo de interesse vs. Brent) e Brent e produzem previsões base de preço. Sobre essa previsão, o método híbrido adiciona duas camadas.

No treinamento dos modelos base, os checkpoints são gerados separadamente para cada combinação ativo-arquitetura e avaliados diretamente na comparação offline. O artigo não realiza seleção \textit{ex post} entre vários checkpoints de uma mesma arquitetura: cada configuração produz um único checkpoint final salvo ao término do treino. As configurações principais são seq\_len igual a 15 para todas as arquiteturas, 140 épocas para LSTM e Autoencoder, 160 épocas para Transformer, taxa de aprendizado de $10^{-3}$ para LSTM e Autoencoder, taxa de $10^{-4}$ para Transformer e, no caso do Autoencoder, peso de reconstrução igual a 0,5.

A primeira camada é o ajuste por sequência recuperada. Dado o conjunto de motivos correntes $M_t$, calculamos a similaridade coseno entre seus embeddings e todos os motivos da biblioteca histórica filtrada. Mantemos os top-$k$ eventos acima do limiar $\tau = 0{,}70$ e calculamos a sequência media ponderada
\[
\bar{q}_h = \frac{\sum_{i \in \mathcal{N}_t} s_i q_{i,h}}{\sum_{i \in \mathcal{N}_t} s_i},
\]
em que $s_i$ e a similaridade do item $i$ e $\mathcal{N}_t$ representa o conjunto de eventos recuperados. A previsão base $\hat{y}^{base}_{t+h}$ e ajustada por
\[
\hat{y}^{seq}_{t+h} = \hat{y}^{base}_{t+h}\left(1 + \alpha \cdot \bar{s} \cdot \frac{\bar{q}_h}{100}\right),
\]
em que $\bar{s}$ e a similaridade media dos eventos selecionados e $\alpha$ e um fator de escala escolhido por busca em grade no conjunto $\{0.1, 0.2, 0.3, 0.4, 0.5, 0.6\}$. Na avaliação offline atual, a recuperação usa top-$k=5$, limiar de similaridade $\tau = 0{,}70$ e aplica o ajuste sequencial no horizonte D0--D4, ainda que a biblioteca armazene sequências até D5 (talvez reduzir a propagação).

A segunda camada é uma correção residual em regime \textit{walk-forward}. Definimos o residual
\[
r_t = y_t - \hat{y}^{seq}_{t},
\]
e treinamos iterativamente um modelo Ridge com atributos derivados apenas do passado observado:
\[
x_t = [r_{t-1}, r_{t-2}, sim_t, event_t, hist_t],
\]
em que $sim_t$ representa a similaridade media do evento corrente, $event_t$ e um indicador binário de evento e $hist_t$ e a quantidade de eventos históricos usados na recuperação. A previsão final é então dada por
\[
\hat{y}^{hyb}_{t} = \hat{y}^{seq}_{t} + g(x_t),
\]
com $g(\cdot)$ estimado apenas a partir da janela passada. Na implementação atual, o ajuste residual usa janela de 15 observações. Essa formulação não deve ser interpretada como identificação causal formal; ela representa um mecanismo de ajuste informado por eventos históricos comparáveis sob restrição temporal estrita (esse trecho está enigmático).

\section{Avaliação Experimental}

\subsection{Protocolo experimental}
A avaliação foi realizada de forma offline, em regime \textit{walk-forward}, sobre os três ativos e as três arquiteturas base. Para cada combinação ativo-modelo, o pipeline gera quatro trilhas: previsão base, ajuste apenas por sequência recuperada (\textit{SeqOnly}), correção apenas residual (\textit{ResidualOnly}) e previsão hibrida completa. Ao todo, o experimento atual produz 12.159 previsões fora da amostra, distribuídas uniformemente entre nove configurações.

As métricas principais são RMSE e MAE, por refletirem erro de magnitude em escala original, que é o alvo central deste estudo de regressão. O RMSE recebe maior destaque por penalizar mais fortemente erros grandes em dias de choque, enquanto o MAE complementa a leitura com uma medida mais direta e menos sensível a outliers. Como métricas complementares, também reportamos \textit{directional accuracy}, percentual de dias em que o hibrido supera o baseline e análises de robustez por ano. A \textit{directional accuracy} é mantida como medida auxiliar porque informa acerto de sinal, mas não captura a magnitude do erro. Além disso, o pipeline gera logs de decisão por ativo e arquitetura, registrando motivos, similaridade media, quantidade de eventos históricos utilizados e datas de referência recuperadas. Esses artefatos reforçam a auditabilidade da camada informacional.

Do ponto de vista metodológico, três cuidados são centrais no protocolo. Primeiro, a recuperação histórica usa apenas eventos anteriores a data prevista, o que reduz vazamento de informação futura. Segundo, o ajuste residual e estimado em regime \textit{walk-forward}, com atributos calculados apenas a partir do passado observado. Terceiro, a avaliação inclui uma ablação explicita para separar o efeito do ajuste sequencial por similaridade do efeito da correção residual.

Com esse desenho, a avaliação permite discutir de forma integrada não apenas a redução agregada de erro frente aos baselines quantitativos, mas também o comportamento do método em dias com e sem evento, a recorrência de sequências semanticamente comparáveis, a robustez entre ativos, arquiteturas e recortes anuais, e os limites da abordagem em cenários de histórico pouco comparável ou sob diferentes limiares de detecção. Assim, os resultados a seguir são interpretados menos como respostas isoladas a perguntas independentes e mais como evidências complementares sobre desempenho, estabilidade, auditabilidade e sensibilidade do método proposto.


\subsection{Resultados principais}

A Tabela~\ref{tab:rmse_principal} resume os resultados de RMSE por ativo e arquitetura. O ganho do método hibrido aparece em todos os nove cenários avaliados. A melhora absoluta varia de 0.5167 a 1.4234 pontos de RMSE, com destaque para EXXO34 com Transformer.

\begin{table}[H]
\centering
\small
\caption{RMSE dos modelos base e híbridos por ativo e arquitetura}
\label{tab:rmse_principal}
\begin{tabular}{l l c c c}
\toprule
\textbf{Ativo} & \textbf{Modelo} & \textbf{RMSE Base} & \textbf{RMSE Hibrido} & \textbf{Ganho} \\
\midrule
PETR4  & LSTM         & 0.9632 & 0.4465 & 0.5167 \\
PRIO3  & LSTM         & 1.5592 & 0.8448 & 0.7144 \\
EXXO34 & LSTM         & 2.0610 & 1.1628 & 0.8982 \\
PETR4  & Autoencoder  & 1.0216 & 0.4487 & 0.5730 \\
PRIO3  & Autoencoder  & 1.6077 & 0.8510 & 0.7567 \\
EXXO34 & Autoencoder  & 2.1269 & 1.1725 & 0.9544 \\
PETR4  & Transformer  & 1.0858 & 0.4624 & 0.6234 \\
PRIO3  & Transformer  & 1.3836 & 0.8439 & 0.5397 \\
EXXO34 & Transformer  & 2.8142 & 1.3908 & 1.4234 \\
\bottomrule
\end{tabular}
\end{table}

Em média, o RMSE cai de 1.6248 para 0.8470, o que corresponde a uma reducao agregada de 47,9\%. O MAE médio cai de 1.2445 para 0.6187, uma redução de 50,3\%. Alem disso, a \textit{directional accuracy} media sobe de 0.4791 para 0.5527, e o método hibrido supera o modelo base em 76,8\% dos dias avaliados. Em termos práticos, isso significa que o ganho não fica restrito a um único ativo nem a uma única arquitetura. Ele aparece de forma consistente em LSTM, Autoencoder e Transformer, bem como em empresas mais maduras e mais novas dentro do recorte analisado.

\subsection{Ablação e interpretação do ganho}

A Tabela~\ref{tab:ablacao} mostra que a camada \textit{SeqOnly}, isoladamente, gera melhora modesta, enquanto a maior parte do ganho vem da combinação entre recuperação semântica e correção residual \textit{walk-forward}. Isso indica que o valor do método não esta apenas em identificar eventos similares, mas em usar essa informação como insumo para uma adaptação residual sensível ao contexto informacional.

\begin{table}[H]
\centering
\caption{Ablação agregada por media de RMSE}
\label{tab:ablacao}
\begin{tabular}{l c}
\toprule
\textbf{Configuração} & \textbf{RMSE médio} \\
\midrule
Modelo base      & 1.6248 \\
SeqOnly          & 1.5789 \\
ResidualOnly     & 0.8881 \\
Hibrido completo & 0.8470 \\
\bottomrule
\end{tabular}
\end{table}

Os números da ablação mostram três mensagens importantes. Primeiro, há algum valor na recuperação semântica pura, mas ela sozinha reduz pouco o erro médio (2,8\%). Segundo, a correção residual orientada por informação de evento e o principal motor de ganho, reduzindo o RMSE medio em 45,3\% frente ao baseline. Terceiro, a combinação das duas etapas ainda acrescenta melhora adicional, levando o ganho agregado a 47,9\%.

\subsection{Resultados por regime e robustez temporal}
Ao separar as observações em dias marcados pela biblioteca de eventos e dias sem marcação, o híbrido melhora em ambos os casos. Em termos médios de MAE, a redução e de 48,4\% nos dias com evento e de 54,6\% nos dias sem evento. Esse resultado sugere que a camada residual aprende regularidades úteis mesmo fora dos dias explicitamente marcados. Ainda assim, o método continua dependente de uma base semântica coerente para capturar os choques mais informacionais.

As análises auxiliares também indicam robustez temporal relevante: 60 de 63 subgrupos ano-ativo-modelo melhoram (o desempenho) em relação ao baseline. As três exceções aparecem em janelas iniciais de 2020 com apenas três observações. Esse comportamento deve ser interpretado à luz da heterogeneidade da amostra. O recorte inclui empresas mais maduras e séries mais longas, como Petrobras, e empresas relativamente mais novas no mercado observado, o que reduz a quantidade de observações disponíveis em janelas muito antigas para alguns ativos. Por isso, os poucos casos sem ganho concentram-se em faixas temporais muito curtas e não indicam degradação sistemática do método.

\subsection{Discussão}
Em conjunto, os resultados sugerem que a biblioteca de eventos recuperados funciona como um mecanismo de memória externa útil para previsão numérica. Seu valor prático, contudo, depende de duas condições: uma base histórica minimamente rica e um procedimento de recuperação sob restrição temporal (?). Isso ajuda a explicar por que a combinação entre recuperação e residual supera tanto o baseline puro quanto o ajuste semântico isolado (tb ficou confuso). O método também preserva uma propriedade importante para uso acadêmico: cada ajuste pode ser rastreado a eventos históricos concretos, o que facilita analise qualitativa, inspeção de falhas e discussão econômica dos resultados.

Como complemento interpretativo, a Tabela~\ref{tab:topicos_impacto} resume os cinco assuntos com maior alteração absoluta de preço no horizonte futuro entre os clusters com pelo menos cinco eventos. O critério usado foi a soma absoluta da trajetória media posterior no horizonte D1--D5 (não seria D0--D4?), preservando o sinal médio observado. Assim, a tabela não mede a polaridade semântica intrínseca (?) do assunto, mas sim o comportamento médio dos preços após eventos daquele tipo.

\begin{table}[H]
\centering
\small
\caption{Assuntos com maior alteração absoluta de preço no horizonte futuro (com maior impacto no preço)}
\label{tab:topicos_impacto}
\begin{tabular}{l c c c}
\toprule
\textbf{Assunto representativo} & \textbf{Ativo} & \textbf{$n$} & \textbf{Impacto D1--D5} \\
\midrule
Preocupações fiscais e politicas no Brasil sobre o ativo & PETR4 & 9 & +6.2527 \\
Redução dos riscos geopolíticos no Oriente Medio & BRENT & 8 & -5.3361 \\
Tensões geopolíticas entre Rússia e Ucrânia & BRENT & 14 & -5.1211 \\
Sanções e restrições ao petróleo russo & BRENT & 12 & -4.3527 \\
Intervenção estatal e politica na gestão da empresa & PETR4 & 7 & -4.0636 \\
\bottomrule
\end{tabular}
\end{table}

\section{Ameaças a Validade}

\subsection{Validade interna}

A principal ameaça a validade interna decorre da dependência de serviços externos para estruturação de notícias (LLM´s ?) e geração de embeddings. Mesmo com prompts controlados, JSON estruturado e \textit{fallback} de busca, ainda existe variação induzida pelo modelo de linguagem e pela disponibilidade de conteúdo recuperado. Além disso, a abordagem não realiza identificação causal formal; portanto, a relação entre notícia, evento recuperado e ajuste preditivo deve ser interpretada como associação informacional útil, e não como prova de causalidade econômica (enigmático, convém siplificar a explicação).

\subsection{Validade de construto}

O construto ``evento relevante'' depende de uma heurística baseada em limiar de retorno diário. Embora a abordagem inclua análise de sensibilidade para 1,5\%, 2,0\% e 2,5\%, ainda há risco de que alguns movimentos importantes fiquem fora da base ou de que certos dias entrem como eventos sem corresponder a um choque informacional claro. De forma semelhante, os clusters semânticos devem ser entendidos como apoio exploratório e não como validação definitiva de categorias econômicas (categorias de notícias?).

\subsection{Validade externa}

O escopo empírico está restrito ao setor de óleo e gás e a três ativos específicos, todos sincronizados com o Brent. Esse desenho foi escolhido por favorecer interpretação econômica e recorrência de choques informacionais, mas limita a generalização imediata para outros setores, mercados ou frequências temporais. A extensão da abordagem para bancos, varejo, tecnologia ou dados intradiários ainda precisa ser demonstrada empiricamente (?).

\subsection{Validade estatística e de reprodução}

Embora o número total de previsões fora da amostra seja expressivo, alguns recortes anuais iniciais possuem poucas observações por causa da maturidade diferente entre os ativos analisados. Empresas mais maduras oferecem janelas históricas mais densas, enquanto ativos relativamente mais novos no recorte observado geram menos pontos em anos iniciais. Essa assimetria afeta sobretudo os recortes temporais mais curtos, justamente onde se concentram os poucos casos sem ganho. Além disso, os artefatos experimentais já incluem logs, tabelas auxiliares e testes de reprodutibilidade, mas ainda não contemplam uma bateria estatística mais ampla, como testes formais de significância, intervalos de confiança por janela ou validação cruzada temporal em múltiplos cortes.

\section{Conclusão}

Este trabalho apresentou um método de previsão híbrida para séries temporais financeiras no qual notícias são convertidas em eventos estruturados e reutilizadas por meio de recuperação semântica de episódios históricos comparáveis. A contribuição central do artigo esta em combinar uma biblioteca histórica de eventos com corte temporal estrito, um ajuste sequencial ponderado por similaridade e uma correção residual \textit{walk-forward} sobre baselines quantitativos padrão.

As evidências quantitativas mostram que essa combinação melhora de forma consistente os três baselines avaliados. O RMSE médio cai 47,9\%, o MAE médio cai 50,3\%, a \textit{directional accuracy} aumenta e o método supera o baseline em mais de três quartos dos dias avaliados. A ablação reforça que o ganho não vem de um único componente isolado, mas da interação entre recuperação semântica e adaptação residual. Em paralelo, a análise exploratória dos clusters indica quais assuntos aparecem associados a trajetórias médias posteriores de maior magnitude, preservando uma camada adicional de auditabilidade.

Como próximos passos, destacam-se quatro direções.(:) A primeira é reconstruir a base semântica por janelas temporais para (visando) reduzir ainda mais o risco de contaminação histórica. A segunda é ampliar a avaliação para outros (ativos,) setores e mercados. A terceira é incorporar testes estatísticos mais fortes (robustos? essa parte não ficou legal) sobre os ganhos observados. A quarta é congelar artefatos não deterministas para reprodução integral do experimento (não compreendi).

\section*{Acknowledgment}
The authors thank CNPq, CAPES, and FAPERJ for partially sponsoring this research.

\bibliographystyle{IEEEtranN}
\small{
 \bibliography{references}
}
	
\end{document}
